\preface

这门课属于是我本科期间学的最好的一门课之一了（百分制成绩99，等级制A$^+$）。

\paragraph{声明}

本课内容基于
\begin{itemize}
	\item Jackson, Classical Electrodynamics.
	\item Griffiths, Introduction to Electrodynamics.
\end{itemize}
其中，除附录 \ref{chap:vector analysis} 矢量分析和\chapref{chap:relativistic electromagnetic}相对论电动力学是参考Griffiths书外，其余均参考Jackson书。

本笔记将Jackson书中的Chap. 2, 3: Boundary-Value Problems in Electrostatics: I, II合并为了\chapref{chap:boundary-value problems}：静电学中的边界条件问题。

本笔记一切公式均采取\textbf{国际单位制}(Système International d'Unités, SI)，没有例外。


\paragraph{补充核对资料}
以下公开资料可用于交叉核对约定、推导和适用条件：
\begin{itemize}
    \item David Tong, \href{https://www.damtp.cam.ac.uk/user/tong/em/el4.pdf}{Electromagnetism and Relativity}, Secs. 5.1--5.3：Lorentz变换、固有时间与四维势。比较张量公式时须先核对度规符号。
    \item David Tong, \href{https://www.damtp.cam.ac.uk/user/tong/em/el5.pdf}{Electromagnetic Radiation}, Secs. 6.2--6.3：偶极辐射功率与Rayleigh散射。
    \item David Tong, \href{https://www.damtp.cam.ac.uk/user/tong/em/el6.pdf}{Electromagnetism in Matter}, Sec. 7.5.4：因果性与Kramers-Kronig关系。
    \item MIT, \href{https://web.mit.edu/6.013_book/www/chapter12/12.5.html}{Electromagnetic Fields and Energy, Sec. 12.5}：复Poynting定理、辐射电阻和散射偏振。
    \item N. Brunner et al., \href{https://doi.org/10.1103/PhysRevLett.93.203902}{Phys. Rev. Lett. 93, 203902 (2004)}：区分群速度与信号速度的光纤实验。
\end{itemize}

\paragraph{绪论}

在Jackson书的Introduction中，着重强调了以下关键词：
\begin{compactitem}
	\item 力(force)与场(field)
	\item 经典(classical)与量子(quantum)
	\item 线性(linear)与非线性(nonlinear)
	\item 数学上的理想化(mathematical idealization)与真实的物理(truly physics)
	\item 时空(spacetime)
	\item 宏观(macroscopic)与微观(microscopic)
\end{compactitem}
